\section{Full Experimental Results}
\label{app:full-results}

% 中文：表~\ref{tab:strict-invalid-rates-all-methods} 报告了五种方法在七个任务、两种推理接口下的严格格式无效率，其中仅 ReaSon 在 DS 推理下存在明显无效输出。
%SO 与 RS 训练在 DS 与 RF 推理下的完整严格 QWK、Macro-F1、Pearson 和 Accuracy 结果分别见表~\ref{tab:full-so-rs-strict-qwk}、
%\ref{tab:full-so-rs-strict-f1} 和~\ref{tab:full-so-rs-secondary-merged}。表~\ref{tab:full-method-task-level} 和~\ref{tab:full-method-averages} 
%进一步在任务层面和跨任务平均上，将 RS、ReaSon、PaIR 和 RePaIR 与 SO 基线进行对比。
Table~\ref{tab:strict-invalid-rates-all-methods} reports the strict-format
invalid rates of five methods across seven tasks and two inference
interfaces, where only ReaSon shows substantial invalidity under DS
inference. The complete strict QWK, Macro-F1, Pearson, and Accuracy results
for SO and RS training under DS and RF inference are given in
Tables~\ref{tab:full-so-rs-strict-qwk}, \ref{tab:full-so-rs-strict-f1},
and~\ref{tab:full-so-rs-secondary-merged}, respectively.
Tables~\ref{tab:full-method-task-level} and~\ref{tab:full-method-averages}
further compare RS, ReaSon, PaIR, and RePaIR against the SO baseline at the
task level and in cross-task averages.

% 中文：表~\ref{tab:full-nll-diagnostics}、\ref{tab:full-rationale-quality-task} 和~\ref{tab:eval-dimensions} 分别展示了 SO 与 RS 的教师强制 NLL 诊断、
%外部大模型对理由质量的评判结果以及评判维度的定义。
Tables~\ref{tab:full-nll-diagnostics}, \ref{tab:full-rationale-quality-task},
and~\ref{tab:eval-dimensions} present the teacher-forced NLL diagnostics for
SO and RS, the external-LLM judgments of rationale quality, and the
definitions of the evaluation dimensions, respectively.

% 中文：表~\ref{tab:full-interface-transitions}、\ref{tab:full-sentence-errors} 和~\ref{tab:full-first-error-position} 
%审计了 DS 到 RF 接口切换的样本级变化、有害事件中理由句的错误类型分布以及首个错误出现的位置。表~\ref{tab:full-correction-aggregate}、
%\ref{tab:full-correction-task} 和~\ref{tab:full-correction-transitions} 报告了理由修正干预的总体恢复结果、按任务分解的恢复率和分数转移分布。
%表~\ref{tab:pair-mix-ablation} 展示了配对损失聚合方式的消融实验：
Tables~\ref{tab:full-interface-transitions}, \ref{tab:full-sentence-errors},
and~\ref{tab:full-first-error-position} audit the per-sample changes after
switching from DS to RF inference, the sentence-level error types in harmful
events, and the position of the first rationale error.
The rationale-correction intervention is reported in
Tables~\ref{tab:full-correction-aggregate},
\ref{tab:full-correction-task}, and~\ref{tab:full-correction-transitions},
covering the aggregate recovery results, per-task recovery rates, and score
transitions.Table~\ref{tab:pair-mix-ablation} shows the ablation of paired-loss
aggregation under matched paired supervision.
